\documentclass[12pt,a4paper]{report}

% ===== PACKAGES =====
\usepackage[utf8]{inputenc}
\usepackage[T1]{fontenc}
\usepackage{times}
\usepackage[margin=1in]{geometry}
\usepackage{graphicx}
\usepackage{setspace}
\usepackage{titlesec}
\usepackage{tocloft}
\usepackage{hyperref}
\usepackage{enumitem}
\usepackage{longtable}
\usepackage{booktabs}
\usepackage{caption}
\usepackage{float}
\usepackage{listings}
\usepackage{xcolor}
\usepackage{fancyhdr}

% ===== SETTINGS =====
\onehalfspacing
\setlength{\parindent}{0pt}
\setlength{\parskip}{6pt}

\hypersetup{
    colorlinks=true,
    linkcolor=black,
    urlcolor=blue,
    citecolor=black
}

\lstset{
    basicstyle=\ttfamily\small,
    backgroundcolor=\color{gray!10},
    frame=single,
    breaklines=true,
    captionpos=b
}

\pagestyle{fancy}
\fancyhf{}
\fancyhead[R]{\thepage}
\fancyhead[L]{\leftmark}
\renewcommand{\headrulewidth}{0.4pt}

% ===== TITLE PAGE =====
\begin{document}

\begin{titlepage}
\centering
\vspace*{1cm}
{\LARGE \textbf{SMART CAMPUS EVENT MANAGEMENT SYSTEM}}\\[0.5cm]
{\LARGE \textbf{(SCEMS)}}\\[2cm]
{\large A Project Report}\\[0.3cm]
{\large Submitted in partial fulfillment of the requirements for the degree of}\\[0.3cm]
{\large \textbf{Bachelor of Technology}}\\[0.3cm]
{\large in}\\[0.3cm]
{\large \textbf{Computer Science \& Engineering}}\\[2cm]
{\large Submitted by}\\[0.3cm]
{\large \textbf{RAMAKRISHNAN S (VTU24465)}}\\[2cm]
{\large Department of Computer Science \& Engineering}\\[0.3cm]
{\large Academic Year 2025--2026}\\
\vfill
\end{titlepage}

% ===== BONAFIDE CERTIFICATE =====
\newpage
\thispagestyle{empty}
\begin{center}
{\LARGE \textbf{BONAFIDE CERTIFICATE}}\\[1.5cm]
\end{center}

This is to certify that the project report titled \textbf{``Smart Campus Event Management System (SCEMS)''} is a bonafide record of work done by \textbf{RAMAKRISHNAN S (VTU24465)} during the Winter Semester of the Academic Year 2025--2026, submitted in partial fulfillment of the requirements for the degree of Bachelor of Technology in Computer Science \& Engineering.

\vspace{3cm}

\begin{tabular}{p{7cm}p{7cm}}
\textbf{Guide} & \textbf{Head of Department} \\[2cm]
\hrulefill & \hrulefill \\
Name \& Designation & Name \& Designation \\
\end{tabular}

\vspace{2cm}
\begin{center}
Date: \hrulefill \hspace{2cm} Place: \hrulefill
\end{center}

% ===== ABSTRACT =====
\newpage
\thispagestyle{empty}
\begin{center}
{\LARGE \textbf{ABSTRACT}}
\end{center}
\addcontentsline{toc}{chapter}{ABSTRACT}
\vspace{1cm}

The management of large-scale campus events traditionally relies on fragmented communication, manual spreadsheets, and physical registrations, leading to significant informational gaps and administrative delays. This report presents the \textbf{Smart Campus Event Management System (SCEMS)}, a comprehensive digital ecosystem designed to centralize and automate the lifecycle of university events.

SCEMS utilizes a modern 3-tier decoupled architecture, integrating a high-performance \textbf{React.js} frontend with a robust \textbf{Spring Boot} backend. The system introduces advanced features such as \textbf{Real-Time QR-based Mobile Payment Synchronization}, an \textbf{AI-powered Support Chatbot}, and \textbf{Interactive Venue Mapping} using Leaflet.js. By leveraging a centralized H2/MySQL database infrastructure, SCEMS provides administrators with real-time analytics and capacity management, while offering students a personalized ``Glassmorphism'' dashboard for seamless discovery and registration.

The implementation results demonstrate a 90\% reduction in manual processing time and a significant increase in student engagement metrics through a streamlined, mobile-responsive user experience. This report details the full technical journey from requirements gathering to the deployment of a state-of-the-art campus solution.

\textbf{Keywords:} Full-Stack Development, React.js, Spring Boot, Event Management, QR Payment Simulation, AI Chatbot, Leaflet Maps, REST API, Glassmorphism UI

% ===== LIST OF FIGURES =====
\newpage
\listoffigures
\addcontentsline{toc}{chapter}{LIST OF FIGURES}

% ===== LIST OF TABLES =====
\newpage
\listoftables
\addcontentsline{toc}{chapter}{LIST OF TABLES}

% ===== LIST OF ACRONYMS =====
\newpage
\begin{center}
{\LARGE \textbf{LIST OF ACRONYMS AND ABBREVIATIONS}}
\end{center}
\addcontentsline{toc}{chapter}{LIST OF ACRONYMS AND ABBREVIATIONS}
\vspace{1cm}

\begin{longtable}{p{3cm} p{10cm}}
\textbf{SCEMS} & Smart Campus Event Management System \\
\textbf{SPA} & Single Page Application \\
\textbf{REST} & Representational State Transfer \\
\textbf{API} & Application Programming Interface \\
\textbf{CRUD} & Create, Read, Update, Delete \\
\textbf{JPA} & Java Persistence API \\
\textbf{ORM} & Object-Relational Mapping \\
\textbf{CORS} & Cross-Origin Resource Sharing \\
\textbf{HMR} & Hot Module Replacement \\
\textbf{RBAC} & Role-Based Access Control \\
\textbf{UPI} & Unified Payments Interface \\
\textbf{QR} & Quick Response \\
\textbf{UI/UX} & User Interface / User Experience \\
\textbf{CSS} & Cascading Style Sheets \\
\textbf{DOM} & Document Object Model \\
\textbf{JWT} & JSON Web Token \\
\textbf{CSV} & Comma-Separated Values \\
\textbf{DFD} & Data Flow Diagram \\
\textbf{NFT} & Non-Fungible Token \\
\textbf{MFA} & Multi-Factor Authentication \\
\end{longtable}

% ===== TABLE OF CONTENTS =====
\newpage
\tableofcontents

% =============================================
% CHAPTER 1: INTRODUCTION
% =============================================
\chapter{INTRODUCTION}
\pagenumbering{arabic}

\section{Introduction}
\linespread{1.5}

In the contemporary academic landscape, university campuses serve as vibrant ecosystems of extracurricular activity, encompassing technical symposiums, hackathons, industrial workshops, cultural festivals, dance competitions, and sporting tournaments. Despite this rich tapestry of events, the administrative processes governing their management remain archaic and fragmented. Event organizers frequently resort to disparate communication channels---physical notice boards, unstructured WhatsApp groups, scattered email threads, and standalone Google Forms---to disseminate information and collect registrations. This fragmentation creates significant data silos, where critical information about event capacity, student participation trends, and departmental engagement metrics remains inaccessible to decision-makers.

The consequences of this disjointed approach are manifold. Students routinely miss opportunities that align with their academic interests because there is no centralized discovery mechanism. Administrators struggle to compile participation data for accreditation reports, often spending days manually reconciling spreadsheets from multiple departments. Payment collection for paid workshops and technical events is handled through informal cash transactions or standalone payment links, creating accountability gaps. Furthermore, there exists no standardized mechanism for collecting post-event feedback, making continuous improvement of campus activities an aspirational goal rather than a data-driven process.

The \textbf{Smart Campus Event Management System (SCEMS)} is conceptualized as a ``Digital Campus Nerve Center'' that addresses these challenges through a comprehensive, full-stack web application. SCEMS is not merely a registration portal; it is a sophisticated platform that bridges the gap between event organizers and the student body by providing end-to-end lifecycle management for campus events. The system leverages modern web technologies---\textbf{React.js} for a dynamic, responsive frontend and \textbf{Spring Boot} for a robust, scalable backend---to deliver an experience that rivals commercial platforms like Eventbrite while being tailored specifically for the academic context.

The platform introduces several innovative features that distinguish it from traditional university portals. A \textbf{Real-Time QR-based Payment Synchronization} mechanism enables cross-device financial simulation, where a student can scan a QR code on their mobile device to complete a payment initiated on their laptop. An \textbf{AI-powered Support Chatbot} provides 24/7 assistance for common queries about registration procedures, venue locations, and event schedules. \textbf{Interactive Venue Mapping} powered by Leaflet.js and OpenStreetMap allows students to visualize exact event locations on an interactive map. These advanced integrations collectively demonstrate the feasibility of building production-grade campus solutions using modern full-stack methodologies.

\section{Aim of the Project}

The primary aim of the Smart Campus Event Management System is to revolutionize the way educational institutions manage, promote, and track campus events by replacing manual, error-prone workflows with an automated, data-driven digital platform. The project pursues the following specific objectives:

\begin{enumerate}
    \item \textbf{Centralized Event Discovery:} To establish a single, authoritative source of truth for every event occurring across all departments of the university. Whether an event is organized by the Computer Science department, the Cultural Committee, or the Sports Council, students should be able to discover it from one unified interface equipped with intelligent search, filtering, and categorization capabilities.

    \item \textbf{Streamlined Multi-Event Registration:} To implement an ``Amazon-style'' shopping cart system where students can browse events, add multiple workshops or competitions to their cart, and complete registration for all selected events in a single, guided checkout process. This eliminates the need for repetitive form-filling across different registration portals.

    \item \textbf{Real-World Financial Simulation:} To demonstrate advanced third-party API integration by building a cross-device payment bridge. The system generates dynamic QR codes using the QR Server API, enabling students to simulate UPI payments by scanning a code on their mobile device while the desktop client polls the backend for real-time status updates.

    \item \textbf{AI-Driven Student Support:} To deploy an intelligent chatbot capable of handling frequently asked questions about event schedules, registration deadlines, venue directions, and cancellation policies. The chatbot employs a keyword-weighted mapping algorithm to provide contextually relevant responses, reducing the burden on event organizers.

    \item \textbf{Data-Driven Administrative Insights:} To empower university administrators with real-time analytics dashboards featuring interactive charts and visualizations. Administrators can monitor registration trends across departments, identify the most popular event categories, track seat utilization rates, and export registration data to Excel format for offline reporting.

    \item \textbf{Interactive Geographical Visualization:} To integrate Leaflet.js-based interactive maps that display the precise geographical location of event venues on the campus, helping students---especially freshers unfamiliar with the campus layout---navigate to their registered events efficiently.
\end{enumerate}

\section{Scope of the Project}

The scope of SCEMS encompasses the complete lifecycle of campus event management, from event creation by administrators to post-event feedback submission by students. The system is designed to serve two primary user roles with distinct functional capabilities.

\subsection{Functional Scope}

\textbf{Student-Facing Capabilities:}
\begin{itemize}
    \item \textbf{Account Management:} Students can authenticate using their institutional credentials. The system supports role-based login, distinguishing between student and administrator access levels.
    \item \textbf{Event Discovery and Search:} A comprehensive browsing interface allows students to explore events organized by category (Technical, Cultural, Sports, Hackathon, Music, Dance, etc.), filter by department, and search using natural-language keywords.
    \item \textbf{Cart-Based Registration:} Students can add multiple events to a persistent shopping cart. Cart state is maintained across browser sessions using \texttt{localStorage}.
    \item \textbf{Simulated Payment Processing:} The checkout flow generates a unique Transaction ID and displays a QR code. The student scans this code on their mobile device, triggering a cross-device synchronization mechanism.
    \item \textbf{Personal Dashboard:} A dedicated dashboard displays the student's registered events, categorized by status: Upcoming, Completed, and Cancelled. Students can unregister from upcoming events.
    \item \textbf{MockMail Inbox:} A simulated email inbox displays professional registration confirmations and receipt notifications.
    \item \textbf{Feedback Submission:} After attending an event, students can rate their experience on a 5-star scale and leave written comments.
\end{itemize}

\textbf{Administrator-Facing Capabilities:}
\begin{itemize}
    \item \textbf{Event CRUD Operations:} Administrators can Create, Read, Update, and Delete events through an intuitive modal-based interface.
    \item \textbf{Real-Time Analytics Dashboard:} Interactive charts powered by Recharts display registration trends (Area Charts) and department-wise participation distribution (Pie Charts).
    \item \textbf{Registration Management:} Administrators can view registrations filtered by event and export the data to Excel/CSV format.
\end{itemize}

\subsection{Technical and Architectural Scope}
The technical scope covers the implementation of a 3-tier decoupled architecture comprising a React.js Single Page Application (frontend), a Spring Boot RESTful API (backend), and an H2 in-memory database (data layer). The system integrates with external third-party APIs including the QR Server API for dynamic QR code generation and OpenStreetMap tile servers via Leaflet.js for interactive venue mapping.

\subsection{Limitations and Out-of-Scope}
The current implementation does not include real payment gateway integration (e.g., Razorpay, Stripe), SMS/Email-based OTP verification, or native mobile application support. These items are identified as future enhancements in Chapter 6 of this report.

\section{Frameworks and Technologies}

The development of SCEMS leverages a carefully selected stack of modern frameworks and technologies, each chosen for its specific strengths in addressing the system's architectural requirements.

\subsection{Frontend Framework: React.js with Vite}
React.js (version 18) serves as the foundation of the presentation layer. Its component-based architecture promotes code reusability, while the virtual DOM ensures high-performance rendering even with large event lists. The Vite build tool was chosen over traditional alternatives like Webpack due to its significantly faster Hot Module Replacement (HMR) during development.

\subsection{UI and Animation Libraries}
The visual identity of SCEMS is built on a ``Glassmorphism'' design language, characterized by semi-transparent backgrounds with CSS \texttt{backdrop-filter} properties. Tailwind CSS provides utility-first styling. Framer Motion handles all page transitions and micro-interactions. Lucide-React provides a cohesive, modern icon set across the interface.

\subsection{Backend Framework: Spring Boot}
Spring Boot (version 3.2) powers the application layer, providing enterprise-grade stability and a rich ecosystem for building RESTful web services. Spring Data JPA with Hibernate serves as the ORM layer. CORS is configured to allow communication between the frontend and backend.

\subsection{Database: H2 with MySQL Compatibility}
The H2 in-memory database is employed for rapid development and testing. Its MySQL compatibility mode ensures that the schema can be seamlessly migrated to a persistent MySQL or PostgreSQL instance for production deployment. The database is automatically seeded with sample data via a \texttt{DataSeeder} component.

\subsection{Third-Party API Integrations}
\begin{itemize}
    \item \textbf{QR Server API (goqr.me):} Generates dynamic QR code images for the payment simulation workflow.
    \item \textbf{OpenStreetMap via Leaflet.js:} Provides interactive, zoomable maps for displaying event venue locations.
\end{itemize}

\begin{table}[H]
\centering
\caption{SCEMS Technology Stack Overview}
\label{tab:tech_stack}
\begin{tabular}{|l|l|l|}
\hline
\textbf{Layer} & \textbf{Technology} & \textbf{Purpose} \\
\hline
Frontend & React.js 18 + Vite & SPA with component architecture \\
\hline
Styling & Tailwind CSS & Utility-first responsive design \\
\hline
Animations & Framer Motion & Page transitions and micro-interactions \\
\hline
Icons & Lucide-React & Modern, consistent iconography \\
\hline
Charts & Recharts & Admin analytics visualization \\
\hline
Maps & Leaflet.js & Interactive venue mapping \\
\hline
Backend & Spring Boot 3.2 & RESTful API and business logic \\
\hline
ORM & Spring Data JPA & Database abstraction layer \\
\hline
Database & H2 (MySQL mode) & Relational data persistence \\
\hline
QR API & goqr.me & Dynamic QR code generation \\
\hline
\end{tabular}
\end{table}

% =============================================
% CHAPTER 2: SURVEY OF SIMILAR APPLICATIONS
% =============================================
\chapter{Survey of Similar Applications in Market}

\section{Generic Event Platforms (Eventbrite, Meetup)}

Eventbrite is widely regarded as the gold standard for public event management. It provides a comprehensive suite of tools for event creation, ticketing, and payment processing, supporting both free and paid events across a global audience. Its strengths include a mature payment gateway integration, mobile-responsive design, and robust analytics for organizers. However, Eventbrite is designed for public-facing events and lacks the institutional specificity required for campus environments. It does not support role-based access control (RBAC) tied to departmental hierarchies, does not integrate with university authentication systems, and its pricing model becomes prohibitive for the volume of small, free events typical of a university calendar.

Meetup.com excels in community building and recurring group activities. Its recommendation engine suggests events based on user interests, and its group management features facilitate ongoing engagement. However, Meetup's architecture is fundamentally oriented toward informal communities rather than structured academic institutions. It lacks support for departmental organization, academic year calendars, and the formal registration workflows (including payment receipts and attendance tracking) that universities require. Furthermore, neither platform provides integrated AI support or cross-device payment simulation capabilities.

\section{Specialized University Portals}

Many universities have developed in-house event management solutions, typically as modules within their existing Enterprise Resource Planning (ERP) systems. These portals are functional but suffer from several critical limitations. First, they are often built on legacy frameworks (e.g., PHP-based systems from the early 2010s) that are not mobile-responsive, resulting in poor usability on smartphones---the primary device used by today's students. Second, their user interfaces are designed for administrative efficiency rather than student engagement, featuring dense tables and form-heavy layouts that discourage casual browsing. Third, these systems operate in isolation from modern web standards, lacking features such as real-time notifications, interactive data visualization, and AI-driven assistance that students have come to expect from consumer applications.

\section{AI-Powered Event Assistants}

A nascent category of event platforms has begun incorporating AI features, primarily in the form of chatbot assistants and recommendation engines. Platforms like Whova and Hopin offer event-specific chatbots for large conferences, but these are typically powered by expensive third-party AI services (e.g., OpenAI GPT) that incur per-query costs. For a university deployment handling thousands of student queries daily, this cost model is unsustainable. SCEMS addresses this by implementing a lightweight, rule-based chatbot that operates entirely on the client side, requiring no external AI API calls while still providing contextually relevant responses to the most common student queries.

\section{Gap Analysis and the SCEMS Advantage}

\begin{table}[H]
\centering
\caption{Comparative Analysis of Event Management Platforms}
\label{tab:gap_analysis}
\begin{tabular}{|l|c|c|c|c|}
\hline
\textbf{Feature} & \textbf{Eventbrite} & \textbf{Meetup} & \textbf{Univ. Portal} & \textbf{SCEMS} \\
\hline
Campus RBAC & No & No & Partial & Yes \\
\hline
Dark Mode / Glassmorphism & No & No & No & Yes \\
\hline
Multi-Event Cart & No & No & No & Yes \\
\hline
Cross-Device QR Payment & No & No & No & Yes \\
\hline
AI Chatbot (Zero Cost) & No & No & No & Yes \\
\hline
Interactive Venue Maps & Partial & Yes & No & Yes \\
\hline
Real-Time Analytics & Yes & Partial & Partial & Yes \\
\hline
Excel Export & Yes & No & Partial & Yes \\
\hline
Mobile Responsive & Yes & Yes & No & Yes \\
\hline
\end{tabular}
\end{table}

The gap analysis reveals that SCEMS uniquely combines the visual sophistication of consumer platforms with the institutional specificity of university portals, while introducing novel features (cross-device QR sync, zero-cost AI chatbot) that are absent from all surveyed alternatives.

\end{document}
